\chapter*{INTRODUCTION GÉNÉRALE}
\markboth{\MakeUppercase{INTRODUCTION GÉNÉRALE}}{}
\addcontentsline{toc}{chapter}{INTRODUCTION GÉNÉRALE}
\adjustmtc
\thispagestyle{MyStyle}

% Guide pour la rédaction de l’Introduction Générale

L’\textbf{Introduction Générale} sert à introduire le projet et à donner une vision globale du travail réalisé. Elle doit inclure les éléments suivants :

\begin{itemize}
    \item \textbf{Contexte :} Présentez le domaine dans lequel s'inscrit votre projet et les problématiques existantes.
    \item \textbf{Problématique :} Définissez la problématique spécifique que vous cherchez à résoudre.
    \item \textbf{Objectifs :} Expliquez les objectifs de votre projet et les bénéfices attendus.
    \item \textbf{Méthodologie :} Mentionnez brièvement les étapes suivies pour mener à bien le projet.
    \item \textbf{Plan du rapport :} Décrivez la structure du document en expliquant le contenu de chaque chapitre.
\end{itemize}

\noindent Cette introduction doit être rédigée de manière claire et concise, en mettant en avant l’importance du projet et son impact.


